\def\level{BTS}  % Classe à droite
\def\course{CIEL 2}
\def\chaptername{CIEL2~--~Accompagnement Personnalisé} % Titre au centre
\def\docpart{Equations différentielles} % Partie du cours
\def\docname{C202} % Nom du fichier sans extension
\pagestyle{cours_FP}
\makeheaderBTS{} % Affiche le bandeau

\begin{multicols}{2}

	\begin{exercice_cours}[Déterminer les solutions]{}
		\begin{enumerate}
			\item $y'-\frac{1}{2}y=0$
			\item $2y'-3y=8y+4y'$
			\item $5y'+3y=0$
		\end{enumerate}
	\end{exercice_cours}

	\begin{exercice_cours}[\quad]{}
		Après avoir déterminé une fonction $f$ de la forme $f(x)=mx\e^{2x}$ avec $m \in \R$ solution particulière de l'équation différentielle $(E):\; y'-2y=2\e^{2x}$, déterminer la solution de $(E)$ telle que $f(0)=1$.
	\end{exercice_cours}

	\begin{exercice_cours}[\quad]{}
A l'aide de Géogébra, déterminer une solution particulière $f$ de l'équation différentielle $(E):\; y'+y=2x\e^{-x}$ sous la forme $f(x)=ax^2\e^{-x}$ avec $a \in \R$. Déterminer la solution de $(E)$ telle que $f(0)=1$.
	\end{exercice_cours}

		\begin{exercice_cours}[\quad]{}
			Résoudre les équations suivantes et déterminer la solution vérifiant $f(x_i)=y_i$:
			\begin{enumerate}
				\item $y'-2y=1$ avec $f(0)=2$
				\item $2y'-y=3y'+y-1$ avec $f(0)=1$
			\end{enumerate}
		\end{exercice_cours}


\end{multicols}

		\begin{exercice_cours}[	On considère l'équation différentielle $(E):\; y'+2y=4\e^{-2x}$.]{}


			\begin{enumerate}
				\item Résoudre l'équation homogène associée à $(E)$
				\item A l'aide de géogébra, déterminer le réel $a$ pour que la fonction $p$ définie par $p(x)=ax\e^{-2x}$ soit une solution particulière de $(E)$
				\item En déduire la solution générale de $(E)$
				\item Déterminer la solution de $(E)$ vérifiant $f(0)=5$.
			\end{enumerate}
		\end{exercice_cours}


		\begin{exercice_cours}[\quad]{}
			Dans un circuit électrique, l'intensité du courant $i(t)$ vérifie l'équation: $(E):\; i'(t)+i(t)=t+1$.

			\begin{enumerate}
				\item Résoudre l'équation homogène associée à $(E)$
				\item A l'aide de géogébra, déterminer une solution particulière de $(E)$ de la forme $i_p(t)=at+b$ avec $a,b \in \R$.
				\item En déduire la solution générale de $(E)$
				\item Déterminer la solution de $(E)$ vérifiant $i(0)=0$.
				\item Déterminer la limite de $i(t)$ en $+\infty$.
			\end{enumerate}

		\end{exercice_cours}